% The larger component: exact proof, no computer rank assumption.
\subsection{35 次元の線形成分と非等次元性}
\label{subsec:vertical-component-v31}

$\Pi=0$、すなわち $c=D=0$ とすると、二つの Poisson 方程式はともに恒等的に零になる。
したがって
\[
 \mathcal L=\{\Pi=0\}
 \simeq S_4^{\oplus3}/xS_3\simeq\mathbb A^{35}
\]
は Poisson 構造のスキームに含まれる。
この包含は、単なる高次元の部分族にとどまらない。

\begin{theorem}[垂直な線形成分]
\label{thm:vertical-component-v31}
$\mathcal L$ は Poisson 構造のスキームの既約成分であり、
その一般点では同スキームは滑らかである。
$\mathcal L$ 上の全 $35$ 次元族は、可換多項式環の次数を保つ Ore 拡大により同時に量子化できる。
\end{theorem}

\begin{proof}
$F^*=(x_1^4,x_2^4,x_3^4)$ を取る。
この点で混合方程式を $p\in S_2^{\oplus3}$ の方向に線形化した写像は
\[
 L_{F^*}(p)
 =F^*(x\cdot p)+3p\cdot F^*-(x\cdot p)\operatorname{div}F^*
 \in S_6
\]
である。この写像が単射であることを係数比較で示す。
$\{i,j,k\}=\{1,2,3\}$ として
\[
 p_i=a_i x_i^2+
      \sum_{j\neq i}b_{ij}x_ix_j+
      \sum_{j\neq i}d_{ij}x_j^2+
      e_i x_jx_k
\]
と書く。
$x_i^6$ の係数は $2a_i$ なので、まず $a_i=0$ である。
$x_i^2x_j^4$ の係数は $-3b_{ij}$ なので $b_{ij}=0$ となる。
$x_ix_j^5$ の係数は $-2d_{ij}+b_{ji}$ なので $d_{ij}=0$ となる。
最後に $x_i^4x_jx_k$ の係数は $-3(e_j+e_k)$ であるから、
$e_1=e_2=e_3=0$ である。
よって $L_{F^*}$ は階数 $18$ を持つ。
この $18$ 個の係数を行とし、上の順に係数を列とした小行列の行列式は
\[
 2^3(-3)^6(-2)^6\det
 \begin{pmatrix}0&-3&-3\\-3&0&-3\\-3&-3&0\end{pmatrix}
 =-20155392\neq0
\]
である。

混合方程式は $p$ に関して線形で、その係数は $F$ に関して線形である。
したがって上の小行列式が非零である $F$ の開近傍では、
混合方程式そのものがスキームとして $p=0$ を強制する。
その開集合上で Poisson 構造のスキームは $\mathcal L$ と一致する。
よって $\mathcal L$ は一般点で滑らかな既約成分である。

量子化には Euler 商の線形切断を一つ選ぶ。
例えば $F_1$ に含まれる $x_1$ で割り切れる単項式をすべて消す切断を取れる。
$F=\sum_iF_i\partial_i$ は可換多項式環 $\mathbb C[x_1,x_2,x_3]$ の導分であるから、
\[
 S_\hbar=\mathbb C[x_1,x_2,x_3][[\hbar]]
 [w;\,-\hbar F],\qquad
 wa=aw-\hbar F(a)
\]
は Ore 拡大である。特に $[x_i,w]=\hbar F_i$ であり、
順序付き単項式 $x_1^{a_1}x_2^{a_2}x_3^{a_3}w^b$ が形式基底をなす。
各 $F_i$ は四次なので関係は重み $(1,1,1,3)$ を保つ。
普遍 $F_{\mathrm{univ}}$ の $35$ 個の係数を完備局所基底の極大イデアルに取り、
式の $\hbar F$ を $F_{\mathrm{univ}}$ に置き換えれば、
$\mathbb C[[u_1,\ldots,u_{35}]]$ 上の平坦な形式変形が得られる。
形式変形の HKR 規約に応じた一様なパラメータの定数倍は、この結論を変えない。
\end{proof}

\begin{corollary}[外側の hull の非等次元性]
\label{cor:non-equidimensional-v31}
Poisson 構造のスキームは、少なくとも次元 $35,19,19$ の三つの異なる既約成分を持つ。
したがって二次 Poisson 錐に同定された外側の miniversal hull は非等次元である。
\end{corollary}

\begin{proof}
$35$ 次元成分は $p=0$ である。
滑らかな三次曲線の $19$ 次元成分では $D=0$、$c\neq0$ であり、
対数型の $19$ 次元成分では一般に $D\neq0$ である。
先の二定理と本定理でそれぞれが既約成分であることを示しているので、互いに異なる。
各成分は斉次錐であり原点を含む。
各成分の斉次素イデアルで割った標準次数付き整域の原点完備化は、
斉次成分の直積として記述できる。二つの非零元の最低次数項の積は零にならないので、
この完備化も整域である。異なる斉次素イデアルの完備化は、元の環への縮約により区別される。
完備化は次元を保つので、三十五・十九・十九次元の相異なる形式的極小成分が残る。
したがって外側の hull は非等次元である。
\end{proof}

この結果は成分の全列挙でも、最大次元が $35$ であるという主張でもない。
残る成分と交差の構造は、特異三次形式とその線形対称性を用いて調べるべき別の問題である。
